\documentclass[a4paper,12pt]{article}
\usepackage[utf8]{inputenc}
\usepackage[T1]{fontenc}
\usepackage{lmodern}
\usepackage{geometry}
\usepackage{setspace}
\usepackage{fancyhdr}
\usepackage{hyperref}
\usepackage{amsmath}
\usepackage{amsfonts}
\usepackage{amssymb}
\usepackage{mathtools}
\usepackage{float}
\usepackage{tabularx}
\usepackage{enumitem}
\usepackage{graphicx}
\usepackage{ragged2e}
\usepackage[spanish]{babel}
\usepackage{booktabs}
\usepackage{longtable}
\usepackage{multirow}
\usepackage{xcolor}
\usepackage{microtype}
\usepackage{caption}
\usepackage{textcomp}
\newcommand{\euro}{\texteuro}
\newcommand{\E}{\mathbb{E}}
\newcommand{\Var}{\mathrm{Var}}
\newcommand{\Cov}{\mathrm{Cov}}
\newcommand{\VaR}{\mathrm{VaR}}
\newcommand{\TVaR}{\mathrm{TVaR}}
\addto\captionsspanish{\renewcommand{\tablename}{Tabla}}
\addto\captionsspanish{\renewcommand{\figurename}{Figura}}
\captionsetup{font=small,labelfont=bf}

% Configuración de márgenes
\geometry{left=2.5cm,right=2.5cm,top=3cm,bottom=3cm}

% Configuración de encabezados y pies de página
\pagestyle{fancy}
\fancyhf{}
\renewcommand{\headrulewidth}{0pt}
\fancyfoot[C]{\thepage}

% Configuración de interlineado
\onehalfspacing
\justifying

% Configuración de hipervínculos
\hypersetup{
    hidelinks,
    colorlinks=false
}

\begin{document}

% Primera página (Título, Resumen y Abstract)
\begin{titlepage}
    \centering
    \vspace*{0.1cm}
    \huge\textbf{Tema 6: Preguntas de repaso}\\
    \vspace{1cm}
    \small{Luiz Alberto Luzuriaga Izquierdo}\\
    \vspace{0.1cm}
    Facultad de ciencias económicas, empresariales y turismo. \\
    \vspace{0.1cm}
    Universidad de Alcalá de Henares. \\
    \vspace{0.1cm}
    Curso 2026-2027\\
    \vspace{0.5cm}

{\bfseries Resumen}

\vspace{0.1cm}
\noindent
\begin{justify}
Respuestas razonadas a las tres preguntas de repaso del tema. Tema 6 (Dependencia no lineal y modelización de colas con cópulas): correlación nula con dependencia de cola, la fórmula estándar de Solvencia II como cópula Gaussiana y la elección de familia para crédito y mercado.\\\\
\vspace{0.1cm}
\noindent
\textbf{Palabras clave}: cópulas, dependencia de cola, Solvencia II, correlación, preguntas de repaso.
\end{justify}
\begin{center}
    {\bfseries Abstract}
\end{center}
\vspace{0.1cm}
\noindent
\begin{justify}
Class work for Topic 6 of Actuarial Mathematics of Stability and Solvency of Insurance Activity: review questions on copulas and tail dependence with actuarial reasoning and numerical evidence.\\\\
\vspace{0.1cm}
\noindent
\textbf{Keywords}: copulas, tail dependence, Solvency II, correlation, review questions.
\end{justify}

\end{titlepage}

% Índice
\newpage\pagenumbering{gobble}
\renewcommand{\contentsname}{Índice}
\newgeometry{left=4cm, right=4cm, top=3cm, bottom=2.5cm}
\tableofcontents

\restoregeometry
\newpage
\numberwithin{equation}{section}
\pagenumbering{arabic}
\setcounter{page}{1}


\section{Tema 6: Dependencia no lineal y modelización de colas con cópulas}
\label{sec:tema6repaso}
\noindent
Respuestas razonadas a las tres preguntas de repaso del tema. Material de referencia: diapositivas del Tema 6 (Prof. Hugo González Riera). Las cifras proceden del taller (Excel, Python y R) con las mismas semillas y datos.

\subsection{Pregunta 1 -- $\rho=0$ pero $\lambda_U>0$: significado y ejemplo de negocio}

\noindent\textbf{Enunciado.} Explica con tus propias palabras qué significa que dos variables tengan $\rho=0$ pero $\lambda_U>0$, y pon un ejemplo de negocio donde esto sería relevante.

\paragraph{Planteamiento.} Son dos afirmaciones sobre objetos distintos: $\rho$ es un promedio lineal sobre toda la distribución conjunta; $\lambda_U$ es el límite de una probabilidad condicionada en la cola. Que ambas sean compatibles no es una paradoja: es la razón de ser de las cópulas.

\subsubsection*{1. Qué dice cada número}
\begin{itemize}[leftmargin=*,itemsep=2pt]
\item $\rho(X,Y)=\operatorname{Cov}(X,Y)/(\sigma_X\sigma_Y)=0$: no hay tendencia \emph{lineal} a que $X$ e $Y$ se muevan juntas \emph{en promedio}. No implica independencia (el ejemplo clásico es $Y=X^2$: dependencia exacta y, en la simulación, $\rho\approx0{,}00$ y $\rho_S\approx0{,}00$).
\item $\lambda_U=\lim_{u\to1^-}P(V>u\mid U>u)>0$: si un riesgo ya está en su peor escenario, la probabilidad de que el otro también lo esté \emph{no se desvanece} al endurecer el escenario. Con independencia, $P(V>u\mid U>u)=1-u\to0$.
\end{itemize}
Pueden coexistir porque el número $\rho$ mezcla todas las regiones del plano: en el centro y en los cuadrantes de signo contrario puede cancelar lo que ocurre en la cola.

\subsubsection*{2. Un ejemplo explícito: la cópula $t$ con $\rho=0$}
En la cópula $t$ con parámetros $(\rho,\nu)$,
\[
\lambda_U=\lambda_L=2\,t_{\nu+1}\!\left(-\sqrt{\frac{(\nu+1)(1-\rho)}{1+\rho}}\right),\qquad
\rho=0\;\Rightarrow\;\lambda=2\,t_{\nu+1}(-\sqrt{\nu+1}).
\]
Con $\nu=4$: $\lambda=2\,t_5(-\sqrt5)=0{,}0756$. Se genera con un factor común: $(Z_1,Z_2)\sim N(0,I)$ independientes y $W\sim\chi^2_\nu/\nu$ común, $T_i=Z_i/\sqrt W$. Cuando $W$ es pequeño, \emph{ambas} coordenadas se inflan a la vez; $\operatorname{Cov}(T_1,T_2)=0$, pero $\operatorname{Cov}(T_1^2,T_2^2)>0$: hay dependencia en las \emph{magnitudes}.

Simulación con $2\cdot10^6$ pares (cópula $t$, $\rho=0$, $\nu=4$):
\begin{center}
\small
\begin{tabular}{lrrr}
\toprule
 & $q=0{,}95$ & $q=0{,}99$ & $q=0{,}999$\\
\midrule
$P(V>q\mid U>q)$ & 0,127 & 0,093 & 0,082\\
Si fueran independientes ($1-q$) & 0,050 & 0,010 & 0,001\\
Razón & $\times2{,}5$ & $\times9$ & $\times82$\\
\bottomrule
\end{tabular}
\end{center}
La correlación de Pearson de las normalizadas es $0{,}002$, la de Spearman $0{,}001$ y $\tau=-0{,}000$: ninguna medida de asociación media ve nada, y sin embargo, a $q=0{,}999$, un extremo arrastra a otro con probabilidad 82 veces mayor que la independiente.

\paragraph{Por qué la correlación cancela.} La cópula $t$ es simétrica: infla por igual los cuadrantes $(+,+)$ y $(+,-)$. Para pérdidas, solo el cuadrante conjunto importa; por eso el negocio necesita cópulas asimétricas (Gumbel, Clayton rotada) cuando se sabe que el contagio es \emph{en la dirección adversa}.

\subsubsection*{3. Efecto sobre el capital}
Con los SCR del Caso 1 (mercado 40, suscripción 30 M\euro) y marginales normales, la independencia ($\rho=0$) da $\sqrt{40^2+30^2}=50{,}00$ M\euro. Con cópula $t$ ($\rho=0$, $\nu=4$) el VaR 99,5\,\% simulado es $54{,}97$ M\euro ($+9{,}9\%$): \textbf{4,97 M\euro de capital} que la fórmula con $\rho=0$ no ve.

\subsubsection*{4. Ejemplo de negocio}
Riesgo catastrófico natural (inundación, granizo) y riesgo de contraparte del reasegurador (o impago de los derivados de cobertura): en tiempos normales sus series no se relacionan ($\rho\approx0$), pero una catástrofe grande \emph{coincide} con rebajas de rating, litigios y retrasos de pago del reaseguro justo cuando más se necesita. Es el mismo mecanismo de \guillemotleft{}riesgo de camino equivocado\guillemotright{} que en banca: la exposición crece cuando la solvencia de la contraparte cae.

\subsubsection*{5. Cuidado con el valor asintótico}
Con $\rho=0$ y $\nu=15$ el límite es $\lambda=0{,}0010$ (casi nulo) y con $\nu=10$, $0{,}0069$. Lo relevante para un SCR 99,5\,\% no es el límite sino la probabilidad condicionada a $q$ finito (aquí $0{,}09$ a $q=0{,}99$). Con datos reales el $q=0{,}99$ se apoya en pocos puntos (en una muestra de $n=1.500$, 15), de modo que el estimador es muy ruidoso (ver Debate~3).

\paragraph{Conclusión.} $\rho=0$ más $\lambda_U>0$ significa \guillemotleft{}sin tendencia lineal media, pero con \emph{acumulación en los extremos}\guillemotright{}. Es el punto ciego de Pearson y el motivo para estimar y vigilar la dependencia de cola.

\subsection{Pregunta 2 -- \textquestiondown{}Por qué la fórmula estándar de Solvencia II es una cópula Gaussiana?}

\noindent\textbf{Enunciado.} \textquestiondown{}Por qué se dice que la fórmula estándar de Solvencia II es, matemáticamente, una cópula Gaussiana? \textquestiondown{}Qué implicación tiene esto para las líneas de negocio con riesgo de contagio extremo?

\paragraph{Planteamiento.} Se demuestra en qué condiciones $SCR=\sqrt{\sum_{ij}\rho_{ij}SCR_iSCR_j}$ \emph{es} el VaR 99,5\,\% de la suma, se comprueba por simulación qué pasa fuera de esas condiciones y se cuantifica lo que cuesta ignorar el contagio.

\subsubsection*{1. La demostración}
Sea $L=(L_1,\dots,L_d)\sim N(0,\Sigma)$ con $\sigma_i=SCR_i/z_{0{,}995}$ y matriz de correlación $P$. La suma $S=\sum_iL_i$ es normal con $\operatorname{Var}(S)=\sum_{ij}\rho_{ij}\sigma_i\sigma_j$, luego
\[
\operatorname{VaR}_{99{,}5\%}(S)=z_{0{,}995}\sqrt{\sum_{ij}\rho_{ij}\sigma_i\sigma_j}=\sqrt{\sum_{ij}\rho_{ij}SCR_iSCR_j}.
\]
Normal multivariante = marginales normales + cópula Gaussiana con matriz $P$. La fórmula estándar \emph{es} ese cálculo, con $P$ prescrita por el regulador.

\paragraph{Generalización.} El resultado vale para cualquier distribución \emph{elíptica} (también la $t$ multivariante): toda combinación lineal $w^\top L$ es una variable de la misma familia escalada por $\sqrt{w^\top\Sigma w}$; si las marginales se calibran en el mismo cuantil, $\operatorname{VaR}(S)=\sqrt{SCR^\top P\,SCR}$. Es decir, lo que la fórmula supone no es solo \guillemotleft{}Gaussiana\guillemotright{} sino \emph{elipticidad} (en particular, ausencia de asimetría y de cola conjunta distinta de la marginal).

\subsubsection*{2. Comprobación por simulación (Caso 1, $2\cdot10^6$ escenarios)}
\begin{center}
\small
\begin{tabular}{lcr}
\toprule
Modelo para $(L_{mkt},L_{uw})$ con $\rho=0{,}25$ & Elíptico & VaR 99,5\,\% (M\euro)\\
\midrule
Fórmula estándar & -- & 55,68\\
(a) Normal bivariante & sí & 55,75\\
(b) $t$ bivariante, $\nu=4$ & sí & 56,02\\
(c) Cópula $t(\nu=4)$ con marginales normales & no & 59,80\\
(d) Cópula Gaussiana, marginales lognormales $\sigma=0{,}8$ & no & 54,38\\
(d') Cópula Gaussiana, marginales lognormales $\sigma=1{,}2$ & no & 53,90\\
(e) Cópula Gaussiana, marginales Pareto $\alpha=3$ & no & 54,36\\
\bottomrule
\end{tabular}
\end{center}
(a) y (b) reproducen la fórmula (la diferencia en (b) es el error de Monte Carlo de un cuantil 99,5\,\% de una $t_4$). En cuanto se sale de la familia elíptica la fórmula deja de ser exacta: con cola de dependencia (c) \textbf{infravalora} ($+7{,}4\%$); con marginales asimétricas (d,e) la \textbf{sobrevalora} ($-2$ a $-3\%$). La fórmula no es ni conservadora ni prudente \guillemotleft{}por construcción\guillemotright{}: depende de la estructura real.

\subsubsection*{3. Implicación para líneas con contagio extremo}
La cópula Gaussiana es asintóticamente independiente ($\lambda=0$ si $\rho<1$), de modo que \textbf{da por hecho} que el beneficio de diversificación del nivel medio (aquí $70-55{,}68=14{,}32$ M\euro, el 20,5\,\%) se mantiene en el percentil 99,5. Con una Clayton rotada de la \emph{misma} $\tau$ de Kendall ($\theta=0{,}383$, $\lambda_U=0{,}16$) el VaR simulado es $62{,}00$ ($+11{,}3\%$) y el beneficio cae a $70-62{,}00=8{,}00$ M\euro: \textbf{se evapora el 44\,\% de la diversificación} precisamente en el escenario que fija el capital. El TVaR 99,5\,\% sube de $62{,}57$ a $70{,}76$ ($+13{,}1\%$), que coincide con el recargo del $13\%$ del Caso~1.

\paragraph{Qué hacer.} (i) Si la entidad usa fórmula estándar, evaluar en el ORSA la desviación del perfil de riesgo respecto a los supuestos del SCR (art.~45.1.c de la Directiva 2009/138/CE) con tablas de sensibilidad a la cola; (ii) para los módulos con contagio conocido (mercado--suscripción de caución o crédito, catastrófico--contraparte) aplicar un modelo interno parcial con cópulas con cola; (iii) justificar ante el supervisor cualquier cópula distinta de la implícita (carga de la prueba en la aseguradora).

\subsection{Pregunta 3 -- Clayton, Gumbel o $t$-Student para crédito y mercado}

\noindent\textbf{Enunciado.} Si tuvieras que elegir entre Clayton, Gumbel y una $t$-Student para modelar la dependencia entre el riesgo de crédito y el de mercado de una cartera de inversión, \textquestiondown{}cuál elegirías y por qué?

\paragraph{Planteamiento.} Se razona sobre las pérdidas (cola superior), se compara qué ofrece cada familia, se comprueba por simulación si los datos permiten distinguirlas y se concluye con una recomendación y su contraste.

\subsubsection*{1. Qué necesita el negocio}
Pérdidas de crédito (ampliación de \emph{spreads}, impagos) y de renta variable: en las crisis financieras se materializan a la vez; en las fases de euforia, una caída conjunta de pérdidas no pone en riesgo la solvencia. El interés está, por tanto, en la \textbf{cola superior de las pérdidas}, que puede ser distinta de la inferior.

\subsubsection*{2. Qué ofrece cada familia (misma $\tau=0{,}5$)}
\begin{center}
\small
\begin{tabular}{lcccp{5.6cm}}
\toprule
Familia & Parámetros & $\lambda_L$ & $\lambda_U$ & Comentario\\
\midrule
Clayton (sobre rentabilidades) & $\theta=2$ & 0,707 & 0 & Cola inferior de rentabilidad = pérdidas conjuntas si se rota; asimétrica\\
Gumbel (sobre pérdidas) & $\theta=2$ & 0 & 0,586 & Cola superior de pérdidas; asimétrica\\
$t$ ($\nu=15$) & $\rho=0{,}707,\nu$ & 0,117 & 0,117 & Cola simétrica y débil\\
$t$ ($\nu=4$) & $\rho=0{,}707,\nu$ & 0,397 & 0,397 & Cola simétrica fuerte\\
\bottomrule
\end{tabular}
\end{center}
Clayton y Gumbel son \emph{intercambiables} en $d>2$ (un único $\theta$) y solo admiten dependencia positiva; la $t$ tiene una matriz $P$ completa (cada par su $\rho$), pero un único $\nu$.

\subsubsection*{3. \textquestiondown{}Los datos distinguen las familias? (simulación, $n=1.500$)}
AIC (menor es mejor) de cada modelo ajustado a muestras generadas con distintas verdades:
\begin{center}
\small
\begin{tabular}{lrrr}
\toprule
 & Verdad: Gumbel $\theta=1{,}8$ & Verdad: Clayton rot.\ $\theta=1{,}2$ & Verdad: Gaussiana $\rho=0{,}68$\\
\midrule
Gaussiana & $-861{,}1$ & $-512{,}8$ & $\mathbf{-888{,}3}$\\
$t$-Student & $-886{,}0$ ($\hat\nu=8{,}2$) & $-534{,}1$ ($\hat\nu=7{,}8$) & $-885{,}3$ ($\hat\nu=129$)\\
Clayton & $-510{,}6$ & $-187{,}2$ & $-644{,}6$\\
Clayton rotada & $-880{,}7$ & $\mathbf{-737{,}2}$ & $-660{,}8$\\
Gumbel & $\mathbf{-944{,}1}$ & $-685{,}4$ & $-795{,}7$\\
Frank & $-803{,}1$ & $-484{,}0$ & $-819{,}2$\\
\bottomrule
\end{tabular}
\end{center}
Con asimetría real y $n=1.500$ el AIC la detecta con claridad: la $t$ intenta imitar la cola con $\hat\nu\approx8$ pero pierde unos 58 puntos frente a la familia correcta; con datos gaussianos la $t$ se degenera ($\hat\nu=129$) y la Gaussiana gana.

\subsubsection*{4. Qué cuesta equivocarse (ejercicio de ajuste, VaR 99,5\,\% de $X+Y$, muestra $t_{15}$)}
Elegir una cola \emph{con la orientación equivocada} es lo más caro: Clayton sin rotar $-13{,}2\%$ (respecto a la $t$ ajustada), Frank (sin cola) $-8{,}4\%$, Gaussiana $-1{,}2\%$, Clayton rotada $+4{,}6\%$, Gumbel $+3{,}9\%$.

\subsubsection*{5. Recomendación}
\begin{enumerate}[leftmargin=*,itemsep=2pt]
\item \textbf{Modelo base: cópula $t$} (matriz $P$ ajustada por pares, $\nu$ estimado), porque crédito y mercado se mueven juntos en ambas direcciones, se generaliza a $d$ riesgos, se calibra con una matriz de correlación coherente con el marco regulatorio y su $\nu$ cuantifica el contagio.
\item \textbf{Alternativa/estrés: Gumbel sobre pérdidas} (o Clayton rotada) si el contraste de asimetría de las colas lo apoya: es más severa a igual $\tau$ y es la prueba de robustez habitual.
\item Decidir con AIC/BIC, bondad de ajuste ($S_n$), $\lambda$ empírico con IC y sensibilidad del capital; \textbf{no} Gaussiana ni Frank como modelo final si hay evidencia de cola; \textbf{nunca} Clayton sin rotar con pérdidas.
\end{enumerate}
\paragraph{Conclusión.} La $t$ es la mejor opción como punto de partida general; si los datos o la economía muestran contagio solo en la dirección adversa, Gumbel (o Clayton rotada) es más fiel; Clayton directa es la cola equivocada para pérdidas.

\subsubsection*{Referencias}
\begin{itemize}[leftmargin=*,itemsep=0pt]
\item Sklar, A. (1959). Fonctions de r\'epartition \`a $n$ dimensions et leurs marges. \emph{Publ. Inst. Statist. Univ. Paris}, 8.
\item Nelsen, R. B. \emph{An Introduction to Copulas}. Springer.
\item Embrechts, P., McNeil, A. y Straumann, D. (2002). Correlation and dependence in risk management: properties and pitfalls.
\item McNeil, A. J., Frey, R. y Embrechts, P. \emph{Quantitative Risk Management}. Princeton University Press.
\item Li, D. X. (2000). On default correlation: a copula function approach. \emph{Journal of Fixed Income}, 9(4).
\item Genest, C., Ghoudi, K. y Rivest, L.-P. (1995). A semiparametric estimation procedure of dependence parameters in multivariate families of distributions. \emph{Biometrika}, 82.
\item Demarta, S. y McNeil, A. J. (2005). The $t$ copula and related copulas. \emph{International Statistical Review}, 73.
\item Directiva 2009/138/CE (Solvencia II) y Reglamento Delegado (UE) 2015/35.
\end{itemize}

\end{document}
